\chapter{Deneyler ve Sonuçlar}
\label{ch:deneyler}

Bu bölümdeki her sayı, üç etiketten biriyle işaretlenmiştir:
\textbf{Ölçüldü (cihazda)} -- Jetson AGX Orin üzerinde gerçek bir koşudan
alınmış ve alıntılanabilir; \textbf{CPU'da, yapısal} -- CUDA'sız bir
geliştirme makinesinde ölçülmüş, bir ilişkinin var olduğunu kanıtlar
ama cihazdaki mutlak maliyeti göstermez; \textbf{TODO} -- araç hazır,
cihazda henüz çalıştırılmamış. Bu üçünü ayırt etmemek, bu projede daha
önce yanlış sonuçlara yol açmıştır (\cref{sec:reprod-tehlike}).

\section{Hesaplama Nerede Gerçekten Yoğunlaşıyor}
\label{sec:flop-dagilimi}

\textbf{Soru.} Görüntü kodlayıcı gerçekten darboğaz mı?

\textbf{Yöntem.} \texttt{tools/analyze\_cost.py} -- \texttt{torch.utils.}
\texttt{flop\_counter.flop\_registry} üzerinden özel bir
\texttt{TorchDispatchMode} ile FLOP sayımı, gerçek kontrol noktasında,
1024×1024.

\begin{yapisal}
\Cref{tab:flop-dagilimi-arkaplan}'daki dağılım (\S\ref{sec:proje-eslesme}):
bellek dikkati \%61,9, görüntü kodlayıcı \%26,8, geri kalan \%11,3.
\end{yapisal}

\textbf{Okuma.} Yalnızca kodlayıcıyı hızlandırmak işin kabaca dörtte
birini karşılamıştır. PyTorch'ta bırakılan bellek dikkati, çoğunluğu
oluşturmaktadır -- \cref{ch:giris}'teki başlangıç varsayımını çürüten
bulgu budur.

\section{Yeniden Yazımın Kayıpsızlığı}

\textbf{Soru.} Dört modül dışa aktarılabilir hale getirilmek için
yeniden yazıldı -- sabit şekilli bellek, karmaşık sayı yerine gerçek
aritmetikli RoPE, toplamsal maskeleme. Bu, modelin hesapladığı şeyi
değiştiriyor mu?

\textbf{Yöntem.} \texttt{tools/analyze\_precision.py --verify-graphs} --
bir klibi yeniden yazılmış modüllerden fp32'de geçirip stok EdgeTAM ile
maskeleri karşılaştırma, gerçek kontrol noktası, 1024×1024.

\begin{olcum}
Tüm 20 karede \textbf{IoU 1,0000}, \texttt{strict=True} ile (orijinal yola
herhangi bir sessiz geri düşüş olsaydı hata fırlatırdı).
\end{olcum}

Bu sonuç 48 birim testiyle desteklenmektedir: gerçek-aritmetikli RoPE'nin
karmaşık RoPE ile eşdeğerliği, doldurulmuş+maskelenmiş belleğin
değişken-uzunluklu bellekle eşdeğerliği, birleştirilmiş bellek
kodlayıcı+Perceiver'ın ayrı ayrı çalıştırılanla eşdeğerliği, özelleştirilmiş
SAM başlığının orijinal \texttt{\_forward\_sam\_heads} ile eşdeğerliği
(nesne mevcutken \emph{ve} kapalıyken), ONNX çıktısının iz (trace)
alınan gruptan farklı bir grup boyutunda bile PyTorch ile eşleşmesi, ve
tüm yamalı takipçinin stok EdgeTAM ile 1 ve 2 nesneyle eşleşmesi.

\section{Motor Bazında Doğruluk ve Hız}
\label{sec:motor-dogruluk-hiz}

\textbf{Soru.} Her fp16 motoru kendi PyTorch modülünü doğru
üretiyor mu, ve ne kadar hızlı?

\textbf{Yöntem.} \texttt{tools/check\_trt\_parity.py} -- fp32 PyTorch
referans, motorun göreli L2 hatası PyTorch'un kendi fp16 otomatik-döküm
(autocast) hatasıyla \emph{yan yana} raporlanır. Üç şekilde zamanlanır:
autocast, TensorRT \texttt{enqueue}, TensorRT CUDA Graph tekrar oynatma.

\begin{olcum}[title={Ölçüldü (cihazda, saatler sabitlenmemiş)}]
\begin{center}
\small
\begin{tabular}{lrrrrr}
  \toprule
  Modül & PyTorch fp16 & TensorRT+Graph & Hızlanma & En kötü L2 & PT fp16'ya oran \\
  \midrule
  Görüntü kodlayıcı & 49,58~ms & 14,30~ms & 3,47× & 2,48e-02 & 2,9× \\
  Bellek dikkati & 14,03~ms & 7,60~ms & 1,85× & 9,10e-04 & 7,9× \\
  Bellek kodlayıcı & 16,63~ms & 2,27~ms & 7,31× & 8,93e-03 & 3,0× \\
  SAM başlığı & 16,70~ms & 1,39~ms & 12,06× & 4,54e-03 & 2,3× \\
  \midrule
  \textbf{Toplam} & \textbf{96,94~ms} & \textbf{25,56~ms} & \textbf{3,79×} & & \\
  \bottomrule
\end{tabular}
\end{center}
\end{olcum}

\textbf{Geçme/kalma kuralı üzerine.} İlk koşu görüntü kodlayıcıyı sabit
bir 2e-2 eşiğine karşı başarısız saymıştı; bu, modelin bir özelliği değil
seçilmiş bir sayıydı ve yanıltıcıydı. Hata, \textbf{derinlik} boyunca
birikmektedir (kodlayıcıda 132 evrişim/doğrusal katman, 2,4e-02 hata; SAM
başlığında 52 katman, 1,0e-03). Bir oran eşiği de tek başına yeterli
değildir: bellek dikkati PyTorch'un fp16'sının 7,9 katında görünür çünkü
autocast orada softmax ve katman normunu fp32'de tutarak alışılmadık
temiz bir taban oluşturur -- mutlak hatası 9e-04 ve önemsizdir. Bu nedenle
bir modül yalnızca hem \texttt{--rel-tol} üzerinde \emph{hem de}
PyTorch'un kendi fp16'sının \texttt{--max-ratio} katı üzerindeyse
başarısız sayılır. Dördü de geçmektedir.

\section{Maskeler Klip Boyunca Değişiyor mu}

\textbf{Soru.} Modül bazlı doğruluk her çağrıda taze rastgele girdi
kullanır. EdgeTAM özyinelemelidir -- ürettiği her maske belleğe yazılıp
sonra geri okunur -- küçük bir modül hatasının biriktiği yüzlerce kare
vardır. Birikiyor mu?

\textbf{Yöntem.} \texttt{tools/compare\_backends.py} -- 500 karelik bir
klibi iki kez izleme (fp32 PyTorch referans, sonra TensorRT fp16), kare
kare IoU ile karşılaştırma.

\begin{olcum}
\begin{center}
\begin{tabular}{ll}
  \toprule
  & \\
  \midrule
  Ortalama IoU & \textbf{0,9993} \\
  En düşük IoU & 0,9973 (kare 405) \\
  0,99 altı & \textbf{500'ün 0'ı} \\
  İlk çeyrek → son çeyrek & 0,9995 → 0,9991 \\
  Eğilim & \textbf{−0,0012 IoU / 1000 kare} \\
  \bottomrule
\end{tabular}
\end{center}
\end{olcum}

\textbf{Okuma.} Önemli olan sayı eğilimdir ve düzdür. Sabit bir ofset ile
yavaş bir sürüklenme bir ortalamada özdeş görünür ama aynı şeyi
söylemez: biri fp16'nın maske sınırındaki yuvarlaması, diğeri
özyinelemeli döngünün onu büyütmesidir. −0,0012/1000 karede döngü
hiçbir şeyi büyütmemektedir.

\section{Uçtan Uca Verim}
\label{sec:uctan-uca-verim}

\textbf{Soru.} Bütün izleme döngüsü 40~ms/kare hedefini tutturuyor mu?

\textbf{Yöntem.} \texttt{tools/benchmark\_tracking.py} -- gerçek
\texttt{propagate()} döngüsü, 500 karelik sentetik klip, stok PyTorch ve
TensorRT arka arkaya aynı kareler üzerinde.

\begin{olcum}[title={Ölçüldü (cihazda, 500 kare, 20 ısınma hariç)}]
\begin{center}
\begin{tabular}{lrrl}
  \toprule
  & ms/kare & FPS & p50 / p90 / p99 \\
  \midrule
  Stok PyTorch (bf16) & 88,30 & 11,32 & 84,35 / 101,29 / 109,13 \\
  \textbf{TensorRT fp16 + CUDA Graph} & \textbf{26,40} & \textbf{37,88} &
    25,01 / 30,84 / \textbf{39,90} \\
  \midrule
  \multicolumn{4}{l}{\textbf{3,34×} hızlanma} \\
  \bottomrule
\end{tabular}
\end{center}
\end{olcum}

\textbf{Okuma.} 40~ms hedefi tutturulmuştur, ve \textbf{p99 de 40~ms'nin
altındadır} -- gerçek zamanlı bir hat için kuyruk, ortalamadan daha çok
önem taşır, çünkü kareyi düşüren kuyruktur.

\section{CUDA Graph'ların Değeri}
\label{sec:cuda-graph}

\textbf{Soru.} TensorRT hedeftir; CUDA Graph'lar yerini hak ediyor mu?

\textbf{Yöntem.} \texttt{tools/benchmark\_tracking.py --cuda-graph-ab} --
aynı arka uç, grafik açık/kapalı, tek işlemde.

\begin{olcum}[title={Ölçüldü (cihazda, 500 kare)}]
\begin{center}
\begin{tabular}{lrrrr}
  \toprule
  & ms/kare & p50 & p99 & p99$-$p50 \\
  \midrule
  Grafik açık & 33,40 & 33,26 & \textbf{37,70} & 4,44 \\
  Grafik kapalı & 35,77 & 37,65 & \textbf{47,18} & 9,53 \\
  \bottomrule
\end{tabular}
\end{center}
\end{olcum}

\textbf{Okuma.} \textbf{CUDA Graph'ların gerekçesi kuyruktur, ortalama
değil.} Grafiksiz p99 47,18~ms'dir -- 40~ms bütçesinin üzerinde; grafikle
37,70~ms'dir -- altında. Titreşim bandı kabaca yarıya iner. Bu tam olarak
grafik tekrar oynatmanın yaptığı şeydir: çağrı başına kernel başlatma
gönderimini kaldırır, ki bu tam olarak başlatma titreşiminin kaynağıdır.

Bu tablodaki \emph{ortalama} fark alıntılanabilir değildir -- aynı
yapılandırma \cref{sec:uctan-uca-verim}'de 26,40~ms, burada 33,40~ms
ölçülmüştür (saatler sabitlenmediği için). Kuyruk argümanı buna
rağmen geçerlidir çünkü iki kol aynı koşuda arka arkaya ölçülmüştür;
ortalama böyle değildir.

\section{fp16, bf16, INT8}
\label{sec:hassasiyet}

\textbf{Soru.} fp16 doğru hassasiyet mi?

\textbf{Yöntem.} \texttt{tools/analyze\_precision.py} -- bir klibi
fp32'de izleyip azaltılmış hassasiyetli etkinleşimlerle yeniden izleme,
maskeleri karşılaştırma. Kasıtlı olarak kötümser: TensorRT'in genellikle
fp32'de bıraktığı normalizasyon ve etkinleşim çıktılarını da yuvarlar,
bu yüzden temiz bir fp16 sonucu iyimser değil güvenlidir.

\begin{olcum}
\begin{center}
\begin{tabular}{lr}
  \toprule
  Hassasiyet & fp32'ye karşı maske IoU \\
  \midrule
  fp16 & \textbf{0,9999} \\
  bf16 & 0,9997 \\
  \bottomrule
\end{tabular}
\end{center}
\end{olcum}

fp16 kazanır ve nedeni yapısaldır: Orin'de bf16 aynı tensör çekirdiklerini
aynı hızda kullanırken 11 bite karşı 8 bit anlamlı basamak taşır ve
TensorRT'in daha az bf16 taktiği vardır.

\textbf{INT8.} Simetrik tensör-başına etkinleşim nicemlemesiyle,
modül bazında:

\begin{itemize}
  \item \textbf{Bellek kodlayıcı}, kırpma yanlılığını biriktiren
        modüldür (0,9397 ortalama IoU, 30 karenin 7'si 0,99 altında).
        Özyinelemeli döngünün \emph{yazma ucudur} -- hatası bankaya
        gider ve geri okunur.
  \item Görüntü kodlayıcı \textbf{yapısal olarak düşman değil,
        kalibrasyona duyarlıdır}: min/maks kalibrasyon onu çökertirken
        yüzdelik dilim kırpma (0,9930) ya da füzyon-modellenmiş
        yerleştirme (0,9926) kurtarmaktadır.
  \item Bütün-model INT8, parça sonuçlarının toplamı değildir.
\end{itemize}

\subsection{INT8 burada bir PyTorch simülasyonudur, TensorRT motoru değil}

\texttt{analyze\_precision.py}'nin \texttt{fake\_quant\_int8*} kancaları
her katmanın \emph{çıktısını} 8 bitlik bir ızgaraya yuvarlayıp geri
float'a çevirir, aksi halde fp32 olan bir ileri geçiş içinde. Yalnızca
etkinleşimler -- ağırlıklara dokunulmaz, hiçbir motor derlenmez.
\texttt{build\_trt\_engines.py}'nin \texttt{PRECISION\_FLAGS}'ı şu anda
yalnızca \texttt{fp16}/\texttt{bf16}/\texttt{fp32} tutmaktadır; bu
depoda bir \texttt{--int8} yolu veya kalibrasyon önbelleği yoktur.
Gerçek INT8 motoruyla ilgili plan \cref{sec:int8-gelecek}'te ele
alınmaktadır.

\section{Ön-İşleme: Derinlemesine Analiz}
\label{sec:onisleme-analizi}

Bu alt bölüm, projede karşılaşılan ve düzeltilen bir ölçüm hatasının tam
kaydıdır -- rapora dahil edilmesinin nedeni, ``kare süresi'' gibi basit
görünen bir ifadenin ne kadar kolay yanlış ölçülebileceğini
göstermesidir.

\subsection{İlk bulgu: ön-işleme küçük değil}

\texttt{pre} ilk kez doğru ölçüldüğünde (kaynak kareyi görselleştirme
için tekrar okumak değil, gerçek kırpma+ölçekleme+normalize), beklenenin
çok üzerinde çıkmıştır:

\begin{olcum}
1024'te \texttt{pre} $\approx$ \textbf{30~ms}, 512'de $\approx$
\textbf{17~ms} -- model süresine (26--33~ms) yakın.
\end{olcum}

Kaynağı çözme değildir; PIL'in yeniden ölçeklemesi ve
\texttt{uint8 / 255.0}'ın NumPy tarafından \textbf{float64}'e
yükseltilmesidir (1024'te kare başına 24~MB dönüşüm + 24~MB
normalizasyon).

\begin{yapisal}[title={CPU'da, yapısal -- 1280×720 JPEG, adım adım}]
\begin{center}
\begin{tabular}{lrr}
  \toprule
  Adım & → 1024² & → 512² \\
  \midrule
  PIL açma + RGB'ye çevirme (çözme) & 4,4~ms & 4,4~ms \\
  PIL yeniden ölçekleme & 18,9~ms & 11,1~ms \\
  \texttt{uint8/255.0} → float64 & 13,6~ms & 0,9~ms \\
  Normalize (float64 üzerinde) & 28,3~ms & 1,3~ms \\
  \midrule
  \textbf{Toplam} & \textbf{$\approx$65~ms} & \textbf{$\approx$17,7~ms} \\
  \bottomrule
\end{tabular}
\end{center}
Aynı işlemler bir dağıtımın yapacağı şekilde: 1024'e yeniden ölçekleme
PIL'de 18,9~ms iken \texttt{cv2.resize} ile 0,6~ms (\textbf{30×});
uint8→float dönüşümü float64'te 13,6~ms iken float32'de 1,9~ms
(\textbf{7×}).
\end{yapisal}

\subsection{Ön-işleme yeniden yazıldı}

Kaynak kare artık doğrudan (JPG önbelleği değil) okunur, kırpma varsa
uygulanır, \texttt{cv2.resize} ile ölçeklenir, float32'de yerinde
(in-place) normalize edilir.

\begin{olcum}[title={Ölçüldü (cihazda, aynı makine ve klip, yalnızca ön-işleme farklı)}]
\begin{center}
\begin{tabular}{lrrr}
  \toprule
  Model girdisi & Önce & Sonra & \\
  \midrule
  1024² & 42,8~ms & \textbf{9,2~ms} & 4,7× \\
  512² & 18,9~ms & \textbf{7,4~ms} & 2,6× \\
  512² (512 merkez kırpmadan) & 5,2~ms & \textbf{5,4~ms} & (zaten yeniden ölçekleme yok) \\
  \bottomrule
\end{tabular}
\end{center}
\textbf{Doğruluk maliyeti:} model girdisi göreli L2 farkı 0,0063 (1024²)
/ 0,0070 (512²); \textbf{maske IoU (eski/yeni), ortalama 0,9993, en düşük
0,9985} -- fp16'nın fp32'ye karşı tüm klipte skorladığı sayının aynısı,
ve görüntü kodlayıcının kendi fp16 hatasından (2,4e-02) $\sim$4× küçük.
\end{olcum}

\subsection{İkinci bulgu: dosya okuma/çözme de bütçenin dışında tutulmalı}

Cihazda \texttt{pre} \textasciitilde40~ms'ye geldi -- ve \textbf{dört
girdi konfigürasyonunun hepsinde} aynı kaldı (1024/512, tam/kırpılmış),
oysa \texttt{inference} çözünürlükle beklendiği gibi ölçeklendi
(26,5 → 8,0~ms). Yeniden ölçeklemeye bağlı hiçbir şey böyle
davranamaz.

\begin{olcum}[title={Ölçüldü (cihazda); 2048×1536 16-bit tek-kanallı TIFF ile doğrulandı}]
\texttt{pre} 40,2~ms, bunun 35,8~ms'si \texttt{cv2.imread}. Aynı dosyada
ayrıştırıldığında: bayt okuma 0,5~ms, \textbf{çözme 34,4~ms}. Yani hariç
tutulan maliyet disk değil, \textbf{görüntü dosyası çözücüsüdür}.
\end{olcum}

Bu ayrım önemlidir çünkü ikisi farklı koşullarda kaybolur: hiçbir
dağıtım diskten TIFF okumaz, ama sıkıştırılmış kare alan bir dağıtım
yine de çözme öder. \texttt{read} artık \texttt{overlay}/kodlama ile
aynı muameleyi görür: ölçülür, raporlanır, bütçeye yazılmaz --
\cref{sec:read-disarida}'de tartışılan koşulla birlikte.

\section{Çözünürlük: 1024'e Karşı 512}
\label{sec:cozunurluk}

\textbf{Soru.} Girdi çözünürlüğünü yarıya indirmek ne kazandırır, neye
mal olur?

\textbf{Yöntem.} \texttt{configs/edgetam\_trt\_512.yaml}, ayrı bir
\texttt{models512/} motor kümesi (motorlar şekle özgüdür). 512, 32×32=1024
uzamsal belirteç üretir -- 1024'teki 4096'ya karşı dört kat azalma.
Karşılaştırma 512 TensorRT'i 512 PyTorch temeline karşı yapılmalıdır;
aksi halde ölçülen şey TensorRT değil çözünürlük değişikliğidir.

\begin{eksik}
\texttt{tools/run\_experiment.py --outdir results/512 --engines models512/
--config configs/edgetam\_trt\_512.yaml --baseline-config
configs/edgetam\_512.yaml --reference-config configs/edgetam\_trt.yaml}
komutuyla doldurulacak:

\begin{center}
\begin{tabular}{lrr}
  \toprule
  & 1024 & 512 \\
  \midrule
  ms/kare (TensorRT) & 26,40 & \\
  FPS & 37,88 & \\
  p99 & 39,90 & \\
  Ortalama IoU (1024 motoruna karşı) & -- & \\
  \bottomrule
\end{tabular}
\end{center}
\end{eksik}

\section{Komut Duyarlılığı}

\textbf{Soru.} Hedef boyutu kare süresini değiştiriyor mu? Nesne sayısı?

\textbf{Yöntem.} \texttt{tools/sweep\_prompt.py} -- her yapılandırma için
ayrı, taze üretilmiş bir klip üzerinde izleme koşusu.

\textbf{Beklenti.} Hedef boyutu düz olmalı: model her kareyi sabit bir
girdiye ölçekler, 12~piksellik bir leke ile 120~piksellik bir leke aynı
şekilli tensör üretir. Nesne sayısı artmalı: nesneler grup (batch)
boyutudur.

\begin{eksik}
\texttt{tools/sweep\_prompt.py --radii 12,30,60,120 --objects 1,2,3}
komutuyla doldurulacak (ms/kare, hedef yarıçapı × nesne sayısı ızgarası).
\end{eksik}

\section{Gerçek Kayıtlar Üzerinde Girdi Konfigürasyonu}
\label{sec:gercek-kayitlar}

\textbf{Soru.} Sentetik klip bunun çalışacağı gerçek veri değildir.
Gerçek kayıtlarda her girdi konfigürasyonu neye mal olur, neyi kaybeder?

\textbf{Yöntem.} \texttt{tools/run\_records.py} -- \texttt{frames/}
altındaki her klasör dört şekilde izlenir:

\begin{table}[H]
  \centering
  \caption{Dört girdi konfigürasyonu.}
  \begin{tabular}{lll}
    \toprule
    Mod & Model & Girdi \\
    \midrule
    \texttt{full1024} & 1024² & tam kare, ölçeklenmiş \\
    \texttt{crop1024} & 1024² & ortalanmış 1024×1024 pencere \\
    \texttt{full512} & 512² & tam kare, ölçeklenmiş \\
    \texttt{crop512} & 512² & ortalanmış 512×512 pencere \\
    \bottomrule
  \end{tabular}
\end{table}

\textbf{Kırpmanın yalnızca üçüncü bir çözünürlük olmadığı.} Dağıtım
kamerası hedefini zaten ortalayıp odaklıyor, yani çerçevenin dışı iki kez
ödenen bir arkaplandır -- bir kez ölçeklenirken, bir kez model üzerinden
geçerken. Tam olarak model girdisine kırpmak ön-işlemeden yeniden
ölçeklemeyi tamamen kaldırır.

\begin{yapisal}[title={CPU'da, yapısal -- aynı klip, aynı 512 model, yalnızca girdi dönüşümü farklı}]
\begin{center}
\begin{tabular}{lrr}
  \toprule
  Girdi & \texttt{pre} medyan & \texttt{post} medyan \\
  \midrule
  Tam 1280×720 kare → 512² & 7,4~ms & 1,5~ms \\
  Ortalanmış 512×512 kırpma, yeniden ölçekleme yok & \textbf{5,4~ms} & 0,7~ms \\
  \bottomrule
\end{tabular}
\end{center}
\end{yapisal}

\textbf{Neye mal olur.} Kırpma, sahnenin 512×512'lik kısmını yerel
çözünürlükte görür; 512 tam ölçekleme, sahnenin tamamını yarı
çözünürlükte görür. Hedef ortalama pencereden çıkarsa \texttt{crop512}
onu tamamen kaybeder.

\begin{eksik}
\texttt{tools/run\_records.py --records frames --out frame\_output}
komutuyla, her kayıt için doldurulacak:

\begin{center}
\begin{tabular}{lrrrr}
  \toprule
  Mod & FPS & \texttt{pre} & \texttt{inference} & \texttt{post} \\
  \midrule
  \texttt{full1024} & & & & \\
  \texttt{crop1024} & & & & \\
  \texttt{full512} & & & & \\
  \texttt{crop512} & & & & \\
  \bottomrule
\end{tabular}
\end{center}

\gorselyertutucu{\texttt{frame\_output/<kayıt>/crop512/tracked.mp4}'ten
alınmış, hedefin kırpma penceresi içinde kaldığını (ya da kaybedildiğini)
gösteren kare(ler)}
\end{eksik}
